\documentclass[11pt,a4paper]{article}

% ── Encodage & langue ──────────────────────────────────────
\usepackage[utf8]{inputenc}
\usepackage[T1]{fontenc}
\usepackage[french]{babel}
\usepackage{lmodern}

% ── Mise en page ───────────────────────────────────────────
\usepackage[a4paper,margin=2.3cm]{geometry}
\usepackage{titlesec}
\usepackage{xcolor}
\usepackage{graphicx}
\usepackage{tabularx}
\usepackage{booktabs}
\usepackage{array}
\usepackage{framed}
\usepackage{fancyhdr}
\usepackage{listings}

% ── Couleurs 3S ────────────────────────────────────────────
\definecolor{bleu3s}{HTML}{1B4F8A}
\definecolor{orange3s}{HTML}{F7941D}
\definecolor{grisclair}{HTML}{F1F5F9}
\definecolor{vertok}{HTML}{15803D}
\definecolor{rougeko}{HTML}{DC2626}

\titleformat{\section}{\Large\bfseries\color{bleu3s}}{\thesection}{0.8em}{}
\titleformat{\subsection}{\large\bfseries\color{bleu3s}}{\thesubsection}{0.7em}{}

\pagestyle{fancy}
\fancyhf{}
\fancyhead[L]{\small\color{bleu3s} 3S TalentMatch --- Module d'Évaluation Technique}
\fancyhead[R]{\small\thepage}
\renewcommand{\headrulewidth}{0.4pt}

\lstset{basicstyle=\ttfamily\small, breaklines=true, frame=single,
        backgroundcolor=\color{grisclair}, framesep=6pt, xleftmargin=8pt}

\begin{document}

% ════════════════ PAGE DE TITRE ════════════════
\begin{center}
  {\Huge\bfseries\color{bleu3s} Module d'Évaluation Technique}\\[6pt]
  {\Large « Reality Gap Score » --- Entretien technique automatisé 100\,\% local}\\[14pt]
  {\large Plateforme 3S TalentMatch --- Projet de Fin d'Études}\\[4pt]
  {\large Youssef GARA --- Juillet 2026}\\[10pt]
  \rule{0.7\textwidth}{1pt}
\end{center}

\vspace{6pt}
\noindent\textbf{Objet du document.} Ce rapport explique le fonctionnement complet du
module d'évaluation : une partie simple (pour tout lecteur), une partie technique
détaillée, et une analyse Internet / confidentialité / RGPD. Il est destiné à la
présentation du travail à l'encadrante et à l'entreprise.

\tableofcontents
\newpage

% ════════════════ PARTIE 1 ════════════════
\section{Explication simple}

\subsection{À quoi sert le module ?}
Un CV \emph{déclare} des compétences, mais rien ne prouve que le candidat les
\emph{maîtrise réellement}. Le module fait passer au candidat un \textbf{entretien
technique automatisé} qui mesure son niveau réel, puis compare ce niveau à ce que son
CV annonce.

\begin{quote}
\textbf{En une phrase :} le CV dit « Expert Python » --- l'évaluation vérifie si c'est vrai.
\end{quote}

\subsection{Le parcours en six étapes}
\begin{enumerate}
  \item \textbf{Le recruteur} sélectionne un ou plusieurs candidats après le matching.
        Il peut ajouter ses propres questions et fixer une date d'ouverture et une date
        limite. Le lancement est \textbf{instantané}.
  \item \textbf{L'IA locale} prépare le questionnaire de l'offre en arrière-plan
        (une seule fois par offre, environ 5 à 8 minutes ; réutilisé ensuite).
  \item \textbf{Le candidat} reçoit un e-mail professionnel : lien personnel, code PIN
        à usage unique, dates de passage. \emph{L'e-mail ne part que lorsque le
        questionnaire est prêt} --- le candidat n'ouvre jamais un test vide.
  \item Il se \textbf{connecte à son compte}, saisit le PIN, lit les \textbf{règles}
        puis passe le test en plein écran : environ 10 QCM adaptatifs (90 secondes
        chacun), 3 questions de raisonnement à rédiger, et les questions du recruteur.
  \item \textbf{L'IA locale note tout} : niveau par compétence, qualité du raisonnement,
        écart entre le CV déclaré et le niveau démontré (« Reality Gap »).
  \item \textbf{Le recruteur consulte} l'onglet « Évaluations » : niveau démontré /10,
        cohérence CV--test, radar « Déclaré vs Démontré », rapport rédigé par l'IA,
        toutes les réponses du candidat, et le score d'intégrité (anti-triche).
\end{enumerate}

\subsection{Ce qui le distingue d'un simple quiz}
\begin{center}
\begin{tabularx}{\textwidth}{>{\bfseries}p{3.2cm}XX}
\toprule
Aspect & Quiz classique & Notre module \\
\midrule
Questions & Banque figée, écrite à la main & Générées par IA à partir de l'offre ---
fonctionne pour tout domaine (IT, finance, santé, droit\ldots) \\
Difficulté & Identique pour tous & Adaptative (IRT) : converge vers le vrai niveau en
$\sim$10 questions \\
Rédaction & Non notée ou par mots-clés & Notée au \emph{sens} par IA sémantique (BGE-M3) \\
Triche & Mêmes questions pour tous & Sous-ensemble différent par candidat, options
mélangées, timer, télémétrie \\
Résultat & Une note brute & Niveau par compétence + écart avec le CV + rapport IA +
réponses vérifiables \\
\bottomrule
\end{tabularx}
\end{center}

% ════════════════ PARTIE 2 ════════════════
\section{Fonctionnement technique détaillé}

\subsection{Architecture}
\begin{lstlisting}
backend/app/
  routes/assessment.py            <- endpoints (20 routes)
  services/assessment/
    llm_client.py                 <- client Ollama LOCAL (aucune API externe)
    question_generator.py         <- generation (paliers + auto-verification)
    question_bank.py              <- pool par offre, arriere-plan, invitations
    cat_engine.py                 <- moteur IRT/CAT (catsim)
    semantic_scorer.py            <- notation semantique (BGE-M3, cosinus)
    reality_gap.py                <- ecart declare CV vs demontre
    report_generator.py           <- rapport IA (Ollama)
  models/assessment.py            <- 4 tables SQLAlchemy
frontend/src/
  components/AssessmentPanel.jsx  <- lancement recruteur
  components/AssessmentsTab.jsx   <- onglet Evaluations (radar, rapport, reponses)
  pages/AssessmentCandidate.jsx   <- passage du test (PIN, regles, QCM, redaction)
  pages/MyAssessments.jsx         <- espace candidat "Mes Evaluations"
\end{lstlisting}

\noindent\textbf{Tables PostgreSQL :}
\begin{center}
\begin{tabularx}{\textwidth}{>{\ttfamily}p{4.6cm}X}
\toprule
\textrm{\textbf{Table}} & \textbf{Contenu} \\
\midrule
assessment\_questions & Pool de QCM par offre (question, 4 options, bonne réponse,
difficulté 1--10, paramètres IRT) \\
open\_questions & Questions de raisonnement + 3 réponses de référence + embeddings \\
assessment\_sessions & Une session par candidat : jeton, PIN, dates, questionnaire tiré,
réponses, theta, scores, signaux d'intégrité \\
reality\_gap\_results & Écart déclaré/démontré par compétence, cohérence CV \\
\bottomrule
\end{tabularx}
\end{center}

\subsection{Génération des questions (100\,\% locale)}
\textbf{Modèle :} Mistral 7B via \textbf{Ollama} (serveur local
\texttt{http://localhost:11434}).

\textbf{Déclenchement :} au premier lancement sur une offre, une tâche d'arrière-plan
(FastAPI \texttt{BackgroundTasks}) génère le pool ($\sim$18 QCM + 6 questions rédigées).
Le recruteur n'attend pas (lancement en 0,2\,s). Les lancements suivants sont immédiats.

\textbf{Entrée du LLM :} \emph{uniquement} le titre du poste et les noms des compétences
requises de l'offre (informations publiques). \textbf{Jamais le CV ni aucune donnée du
candidat.}

\textbf{Trois mécanismes de qualité :}
\begin{enumerate}
  \item \textbf{Paliers de difficulté} : deux appels séparés --- moitié « fondamentaux »
        (difficulté 2--5), moitié « difficile/expert » (7--10 : pièges, cas limites).
  \item \textbf{Distracteurs plausibles imposés} : quatre options proches, de même forme ;
        les mauvaises options sont des erreurs fréquentes du métier $\Rightarrow$ deviner
        au hasard ne donne aucun indice.
  \item \textbf{Auto-vérification} : le LLM répond lui-même à chaque QCM généré
        \emph{sans voir} la réponse marquée ; en cas de désaccord, la question est rejetée.
\end{enumerate}

\subsection{Test adaptatif (IRT / CAT)}
\textbf{Théorie :} \emph{Item Response Theory}, le standard des tests standardisés
(GMAT, TOEFL). Librairie : \texttt{catsim}.

\begin{itemize}
  \item Chaque question possède des paramètres IRT : difficulté $b$ (mappée de 1--10 vers
        $[-3, +3]$), discrimination $a$, pseudo-chance $c = 0{,}25$ (4 options).
  \item Le niveau du candidat est la variable $\theta$ (theta), initialisée à 0.
  \item À chaque réponse, $\theta$ est ré-estimé par maximum de vraisemblance, puis la
        question suivante est choisie parmi les plus \emph{informatives} au niveau
        $\theta$ courant (sélecteur « randomesque » : tirage parmi les 5 meilleures,
        donc deux candidats de même niveau ne voient pas les mêmes questions).
  \item Le test s'arrête vers 10 questions ; $\theta$ est converti en niveau 0--10,
        avec un détail par compétence.
\end{itemize}

\textbf{Pourquoi l'IRT plutôt qu'un pourcentage ?} Un pourcentage dépend des questions
posées (8/10 sur des questions faciles $\neq$ 8/10 sur des dures). $\theta$ estime le
niveau indépendamment des questions : test court \emph{et} comparaison équitable.

\subsection{Notation des réponses rédigées (sémantique)}
\textbf{Modèle :} BGE-M3 --- le même modèle d'embeddings que le moteur de matching
(instance partagée).

Chaque question rédigée possède trois réponses de référence (\emph{faible},
\emph{correcte}, \emph{experte}) dont les vecteurs sont pré-calculés. À la soumission :
\begin{enumerate}
  \item la réponse du candidat est encodée en vecteur (localement) ;
  \item similarité cosinus avec chacune des trois références ;
  \item score 0--100 par interpolation \emph{softmax} vers les ancres 25 / 62 / 95 ;
  \item pénalité $\times 0{,}6$ si la réponse fait moins de 30 mots.
\end{enumerate}
On note donc le \textbf{sens} de la réponse, pas des mots-clés. Validation
expérimentale : experte $\approx 80$ $>$ correcte $\approx 41$ $>$ faible $\approx 16$.

Les \textbf{questions du recruteur} sont posées dans le même flux mais ne sont pas
notées par l'IA : le recruteur les juge lui-même (souvent des questions de motivation
sans « bonne réponse »).

\subsection{Le « Reality Gap » (écart CV -- démontré)}
\begin{enumerate}
  \item \textbf{Niveau déclaré} (0--10) dérivé du CV parsé, par compétence, avec une
        formule transparente : base 5 si la compétence est listée ; $+0{,}4$/année
        d'expérience indiquée (plafond $+3$) ; $+2$ si utilisée dans une expérience
        réelle ; bonus séniorité ($+1$ confirmé, $+2$ senior). Absente du CV : 0.
  \item \textbf{Écart} : $gap = |déclaré - démontré| / 10$ par compétence, moyenne
        pondérée par l'importance dans l'offre (requise $=1{,}0$ ; appréciée $=0{,}5$).
  \item \textbf{Cohérence CV} $= 100 - gap \times 100$, avec les seuils :
        $\geq 85$ « CV cohérent » ; 70--85 « à vérifier » ; $< 70$ « CV $\neq$ niveau
        démontré ».
  \item \textbf{Score final ajusté} $= score_{matching} \times (0{,}7 + 0{,}3 \times
        cohérence/100)$.
\end{enumerate}
L'interface affiche un \textbf{radar superposé « Déclaré (CV) vs Démontré (test) »} par
compétence : le recruteur voit d'un coup d'\oe il les compétences surévaluées.

\subsection{Le rapport IA}
Principe clé : \textbf{le code calcule tous les chiffres} ($\theta$, niveaux, scores,
gap --- aucun risque d'hallucination) ; \textbf{l'IA locale rédige l'analyse
qualitative} : synthèse, niveau technique, qualité du raisonnement, cohérence CV,
verdict (\textsc{recruter} / \textsc{à approfondir} / \textsc{rejeter}) justifié.
Le recruteur peut tout vérifier : bouton « Voir les réponses » (chaque QCM avec le
choix du candidat vs la bonne réponse, chaque rédaction avec son score).

\subsection{Sécurité et anti-triche}
\textbf{Contrôle d'accès --- trois verrous successifs :}
\begin{enumerate}
  \item \textbf{Compte} : le candidat doit être connecté à \emph{son} compte (celui lié
        au CV) ; un administrateur connecté peut aussi ouvrir le lien (test/aperçu).
  \item \textbf{PIN} : code à 6 chiffres, à usage unique, envoyé par e-mail, exigé au
        démarrage et revalidé côté serveur à chaque réponse.
  \item \textbf{Fenêtre de dates} : avant l'ouverture ou après la limite, l'accès est
        refusé.
\end{enumerate}

\textbf{Anti-triche --- traçabilité, pas de sanction automatique :}
\begin{center}
\begin{tabularx}{\textwidth}{p{5.2cm}X}
\toprule
\textbf{Risque} & \textbf{Contre-mesure} \\
\midrule
Partage des questions & Sous-ensemble aléatoire par candidat + options mélangées
(permutation déterministe par session, remappée côté serveur) \\
Recherche externe pendant un QCM & Timer 90\,s par question (temps écoulé = compté faux) ;
temps de réponse enregistré (QCM difficile réussi en moins de 5\,s = signalé) \\
Copier-coller (ChatGPT\ldots) & Bloqué sur les questions et les réponses ; tentatives comptées \\
Texte injecté sans taper & Télémétrie clavier : frappes comptées vs longueur du texte \\
Sortie de l'écran du test & Plein écran obligatoire (annoncé et accepté sur l'écran des
règles) ; chaque sortie et changement d'onglet est enregistré \\
\bottomrule
\end{tabularx}
\end{center}

Ces signaux alimentent un \textbf{score d'intégrité (0--100)} présenté au recruteur avec
la liste des incidents. \textbf{Choix de conception assumé :} le système ne bloque
jamais automatiquement le candidat (un rafraîchissement ou la touche Échap ne prouvent
pas la triche) --- il \emph{trace et signale}, la décision reste \emph{humaine}. C'est
le modèle des outils professionnels de surveillance d'examen, et la position la plus
défendable juridiquement.

% ════════════════ PARTIE 3 ════════════════
\section{Internet, confidentialité et RGPD}

\subsection{Qu'est-ce qui utilise Internet ? (audit du code)}
\begin{center}
\begin{tabularx}{\textwidth}{Xp{4.6cm}>{\centering\arraybackslash}p{2.2cm}}
\toprule
\textbf{Composant} & \textbf{Où il s'exécute} & \textbf{Internet ?} \\
\midrule
Génération des questions (Mistral 7B) & Ollama, sur la machine & \textcolor{vertok}{\textbf{Non}} \\
Notation sémantique (BGE-M3) & En mémoire, sur la machine & \textcolor{vertok}{\textbf{Non}} \\
Test adaptatif (catsim) & Bibliothèque Python locale & \textcolor{vertok}{\textbf{Non}} \\
Rapport IA & Ollama, sur la machine & \textcolor{vertok}{\textbf{Non}} \\
Base de données & PostgreSQL local & \textcolor{vertok}{\textbf{Non}} \\
Envoi des e-mails d'invitation & SMTP & \textcolor{rougeko}{\textbf{Oui}} \emph{(seul flux)} \\
\bottomrule
\end{tabularx}
\end{center}

Le \textbf{seul flux sortant} est l'e-mail d'invitation : adresse du candidat, prénom,
titre du poste, lien, PIN, dates. \textbf{Aucun contenu de CV, aucune réponse, aucun
résultat ne quitte jamais la machine.}

\begin{quote}
\emph{Recommandation production :} remplacer le SMTP actuel (Gmail) par le serveur mail
interne de l'entreprise --- une variable d'environnement à changer, aucun code.
\end{quote}

À noter : une première version du module d'entretien reposait sur une API cloud (Groq) ;
elle a été \textbf{entièrement supprimée} du projet (code, tables, clé API) pour
respecter l'interdiction des API externes. Preuve : \texttt{grep -ri groq backend/}
$\rightarrow$ 0 résultat.

\subsection{Analyse RGPD}
\begin{center}
\begin{tabularx}{\textwidth}{p{4.6cm}X}
\toprule
\textbf{Principe RGPD} & \textbf{Réponse du module} \\
\midrule
Minimisation des données & Le LLM ne reçoit que le titre du poste et les compétences de
l'offre --- jamais le CV ni l'identité. \\
Pas de transfert hors machine & Tout le traitement IA est local (Ollama + BGE-M3).
Aucune donnée candidat n'est envoyée à un tiers. \\
Information préalable & Écran des règles avant le test : le candidat est informé
(plein écran, enregistrement des sorties, analyse du clavier, timer) et clique pour
accepter avant de commencer. \\
Proportionnalité de la surveillance & Pas de webcam, pas de micro, pas de capture
d'écran --- uniquement des signaux techniques, et aucune décision automatique. \\
Décision non automatisée (art.~22) & Le système \emph{assiste} : le recruteur consulte
les réponses brutes et décide seul. \\
Droit d'accès / effacement & Réponses consultables ; suppression d'une évaluation en un
clic (cascade session + résultats). \\
Sécurité d'accès & Triple verrou (compte + PIN à usage unique + fenêtre de dates) ;
endpoints recruteur protégés par authentification et rôle. \\
\bottomrule
\end{tabularx}
\end{center}

\noindent\textbf{À formaliser côté entreprise avant production :} durée de conservation
des sessions/résultats (ex.\ purge à 12 mois) ; mention de l'évaluation dans la
politique de confidentialité ; SMTP interne.

\subsection{Questions probables et réponses}
\begin{description}
  \item[« Pourquoi l'IRT plutôt qu'un simple score ? »] Le niveau estimé est indépendant
        des questions posées : test plus court ($\sim$10 questions), équitable entre
        candidats ; c'est la méthode des tests standardisés.
  \item[« Comment garantir la qualité des questions générées ? »] Paliers de difficulté
        imposés, distracteurs plausibles obligatoires, auto-vérification par le modèle ;
        en dernier ressort le recruteur voit toutes les questions ; et le modèle local
        est remplaçable par un plus performant (une variable d'environnement).
  \item[« Et si le candidat triche ? »] La triche est rendue difficile (questions
        uniques, options mélangées, timer, copier-coller bloqué), détectable (télémétrie
        clavier, temps de réponse, sorties d'écran) et visible (score d'intégrité).
        La sanction reste humaine --- choix éthique et juridique.
  \item[« Quelles données personnelles, et où ? »] CV et réponses : PostgreSQL local.
        Traitement IA : local. Seul envoi externe : l'e-mail d'invitation.
  \item[« Quelle est la contribution originale ? »] Le \textbf{Reality Gap Score} :
        croiser automatiquement le \emph{déclaré} (CV parsé par le pipeline NLP) et le
        \emph{démontré} (test IRT + notation sémantique) pour produire un indice de
        cohérence du CV, avec radar par compétence et score de matching ajusté.
\end{description}

\subsection{Limites (honnêteté scientifique)}
\begin{enumerate}
  \item \textbf{Qualité du modèle local} : Mistral 7B produit occasionnellement une
        question discutable malgré les filtres ($\sim$1/16 observé). Mitigation :
        vérification humaine possible ; modèle remplaçable sans changer le code.
  \item \textbf{Première génération lente} ($\sim$5--8 min par offre) --- mitigée par
        l'arrière-plan, l'invitation différée et le cache par offre.
  \item \textbf{Niveau déclaré dérivé} : les CV n'indiquent presque jamais un niveau
        chiffré (1\,\% des cas mesurés) ; la formule de dérivation est transparente mais
        reste une heuristique.
  \item \textbf{Anti-triche web} : un second appareil ou une tierce personne restent
        indétectables par un navigateur --- d'où le score d'intégrité et la
        contre-vérification humaine en entretien final.
\end{enumerate}

\vspace{14pt}
\begin{center}
\rule{0.5\textwidth}{0.4pt}\\[6pt]
\small Document généré le 11 juillet 2026 --- branche \texttt{feature/reality-gap-score}.
\end{center}

\end{document}
